% from ega3-1-fr.tex
\oldpage[0\textsubscript{III}]{16}
admet un point générique $y$, il existe $\alpha$ tel que $y\in E_\alpha$,
donc $Y=\overline{\{y\}}\subset\overline{E_\alpha}$, et la condition
$2^\circ$ est conséquence de $1^\circ$.

\begin{proposition}[9.2.5]
\label{0.9.2.5-fr}
Soient $X$ un espace noethérien, $x$ un point de $X$, $E$ une partie
constructible de $X$. Pour que $E$ soit un voisinage de $x$, il faut et il
suffit que pour toute partie fermée irréductible $Y$ de $X$ contenant $x$,
$E\cap Y$ soit dense dans $Y$ (s'il existe un point générique $y$ de $Y$,
cela signifie aussi (\hyperref[0.9.2.2-fr]{9.2.2}) que $y\in E$).
\end{proposition}

La condition est évidemment nécessaire~; prouvons qu'elle est suffisante.
Appliquant le principe de récurrence noethérienne à l'ensemble $\mathfrak{M}$
des parties fermées $Y$ de $X$ contenant $x$ et telles que $E\cap Y$ soit un
voisinage de $x$ dans $Y$, on peut supposer que toute partie fermée $Y\ne X$
de $X$ contenant $x$ appartient à $\mathfrak{M}$. Si $X$ n'est pas
irréductible, chacune des composantes irréductibles $X_i$ de $X$ contenant
$x$ est distincte de $X$, donc $E\cap X_i$ est un voisinage de $x$ par
rapport à $X_i$~; par suite, $E$ est un voisinage de $x$ dans la réunion des
composantes irréductibles de $X$ contenant $x$, et comme cette réunion est un
voisinage de $x$ dans $X$, il en est de même de $E$. Si $X$ est
irréductible, $E$ est dense dans $X$ par hypothèse, donc contient une partie
ouverte non vide $U$ de $X$ (\hyperref[0.9.2.2-fr]{9.2.2})~; la proposition
est alors évidente si $x\in U$~; sinon, $x$ est par hypothèse intérieur à
$E\cap(X\setminus U)$ par rapport à $X\setminus U$, donc l'adhérence dans
$X$ de $X\setminus E$ ne contient pas $x$, et le complémentaire de cette
adhérence est un voisinage de $x$ contenu dans $E$, ce qui achève la
démonstration.

\begin{corollary}[9.2.6]
\label{0.9.2.6-fr}
Soient $X$ un espace noethérien, $E$ une partie de $X$. Pour que $E$ soit un
ensemble ouvert dans $X$, il faut et il suffit que pour toute partie fermée
irréductible $Y$ de $X$ rencontrant $E$, $E\cap Y$ contienne une partie
ouverte non vide de $Y$.
\end{corollary}

La condition est évidemment nécessaire~; inversement, si elle est vérifiée,
elle implique que $E$ est constructible en vertu de
(\hyperref[0.9.2.3-fr]{9.2.3}). En outre,
(\hyperref[0.9.2.5-fr]{9.2.5}) montre que $E$ est alors voisinage de chacun
de ses points, d'où la conclusion.

\subsection{Fonctions constructibles}
\label{subsection:0.9.3-fr}

\begin{definition}[9.3.1]
\label{0.9.3.1-fr}
Soit $h$ une application d'un espace topologique $X$ dans un ensemble $T$.
On dit que $h$ est \emph{constructible} si $h^{-1}(t)$ est constructible pour
tout $t\in T$, et vide sauf pour un nombre fini de valeurs de $t$~; pour
toute partie $S$ de $T$, $h^{-1}(S)$ est alors constructible. On dit que $h$
est \emph{localement constructible} si tout $x\in X$ possède un voisinage
ouvert $V$ tel que $h|V$ soit constructible.
\end{definition}

Toute fonction constructible est localement constructible~; la réciproque
est vraie quand $X$ est quasi-compact et admet une base formée d'ensembles
ouverts rétrocompacts dans $X$ (en particulier quand $X$ est noethérien).

\begin{proposition}[9.3.2]
\label{0.9.3.2-fr}
Soit $h$ une application d'un espace noethérien $X$ dans un ensemble $T$.
Pour que $h$ soit constructible, il faut et il suffit que pour toute partie
fermée irréductible $Y$ de $X$, il existe une partie non vide $U$ de $Y$,
ouverte par rapport à $Y$, et dans laquelle $h$ soit constante.
\end{proposition}

La condition est nécessaire~: en effet, par hypothèse, $h$ ne prend dans $Y$
qu'un nombre fini de valeurs $t_i$, et chacun des ensembles
$h^{-1}(t_i)\cap Y$ est constructible dans $Y$
(\hyperref[0.9.2.1-fr]{9.2.1})~; comme ils ne peuvent être tous des parties
rares de l'espace $Y$, un d'eux au moins contient un ensemble ouvert non vide
(\hyperref[0.9.2.3-fr]{9.2.3}).

